\section{Modelo do Sistema}
\label{sec:modelo}

O sistema considera um modelo parcialmente síncrono. Cada chamada RPC possui
um deadline conhecido, mas o escalonamento dos processos e os atrasos da rede
podem variar. Os canais utilizam TCP por meio do gRPC e não há garantia de
entrega quando o processo de destino está indisponível.

As falhas consideradas seguem o modelo de \textit{crash-recovery}: um nó que
falha interrompe abruptamente suas respostas às chamadas e, ao ser restaurado,
reinicia como um novo processo com sua configuração inicial (sem persistência
de estado em disco). O modelo não contempla falhas bizantinas nem partições de
rede que isolem conjuntos disjuntos de nós.

Os $N$ nós que compõem o sistema ($N \ge 3$) são processos independentes,
mesmo quando compartilham a mesma máquina física por meio do Docker Compose. O
maior identificador entre os nós ativos possui prioridade para coordenador. O
estado local inclui o identificador da eleição corrente, o coordenador
conhecido e a indicação de eleição em andamento.

As propriedades esperadas são:
\begin{itemize}
    \item \textit{segurança}: somente o maior nó ativo se declara coordenador
    em uma execução sem mensagens antigas atrasadas;
    \item \textit{vivacidade}: após uma falha permanente, os nós restantes
    concluem uma nova eleição;
    \item \textit{detecção}: um atraso isolado não inicia eleição, pois são
    exigidas duas falhas consecutivas de heartbeat.
\end{itemize}
